\scanpage{45}\mspage{44}
\uline{Corollaire} Si $\deleted{\unsure{X}}Y\ll X$, alors \deleted{la relation} toute ombre ouverte d'une \deleted{\unsure{strate de}} \inserted{strate \unsure{eff.|aff.}} $Y$ est contenue \unsure{dans|ds} l'ombre ouverte d'une strate de $X$. Plus généralement et par suite, pour \unsure{tte} figure $F$ et \deleted{\unsure{t}} \deleted{sous-figure} \inserted{raffinement} $F'$ de $F$, \unsure{toute} \deleted{strate} \deleted{\unsure{de} ombre ouverte d'une} strate ouverte de $\unsure{F'|F}$ est contenue dans une (unique) strate ouverte de $F$.

\deleted{\uline{\unsure{Cond}}. \unsure{8} \unsure{Si} \unsure{/}}\enspace$\backslash$

Je viens de m'apercevoir que dans le cas du pseudo disque, \deleted{\unsure{une}} deux lieux disjoints \deleted{\unsure{cq}} \inserted{$x,y$} de l'intérieur (lesquels définissent donc deux ``raffinements'' dudit pseudo-disque), qu'il existe une ``Subdivision'' du disque dont ils soient des \unsure{noeuds|nœuds}. [\uline{NB}. Les \uline{noeuds} d'une figure \inserted{$\deleted{\unsure{F}}F$} sont les multistrates \inserted{de $\deleted{\unsure{X|F}}\unsure{F}$} réduites à un lieu, i.e. les éléments de $F\cap\unsure{\mathcal L|L}$] Donc \deleted{\unsure{la no}} il n'y a pas lieu de considérer $\{x\unsure{,|.}y\}$ comme un ``raffinement'', vu qu'il ne s'étend pas en une subdivision. Donc il est inapproprié
